% ============================================================
% Introduzione al cheat sheet
% ============================================================
\section*{Come è fatto questo cheat sheet}
\gdef\currentchapter{Introduzione}

Il cheat sheet raccoglie in poche pagine la notazione UML e i design pattern del corso, e accompagna la dispensa
(capitoli 6, 7, 8, 9 e 15). Ha due parti.
\begin{itemize}[itemsep=2pt]
  \item \textbf{Parte I, i diagrammi UML.} Per ogni diagramma le \textbf{voci} della notazione, una per riga: il
        numero e il nome, il disegno, il significato e, quando c'è, l'equivalente in Java. Per il diagramma delle
        classi ci sono poi i \textbf{casi «dal codice al diagramma»}: a sinistra un pezzo di codice Java, a destra il
        diagramma che gli corrisponde e la regola per passare dall'uno all'altro. I casi coprono tutto ciò che si
        può trovare in una classe Java: campi semplici e costanti, riferimenti ad altre classi, \textbf{array},
        \textbf{liste}, insiemi, mappe e matrici, ereditarietà e interfacce, enumerazioni, classi generiche e
        annidate, dipendenze, tutto-parti. Gli ultimi argomenti traducono in codice anche i diagrammi di sequenza,
        di stato e di attività, e una tabella aiuta a scegliere il diagramma giusto.
  \item \textbf{Parte II, riconoscere i design pattern.} Si apre con un metodo per riconoscere un pattern dal testo
        di un esercizio, dal codice o dal diagramma, e con la tabella delle «firme» che distinguono i pattern che si
        somigliano. Poi una sezione per ogni pattern: la \textbf{carta d'identità} (intento, indizi nel testo, nel
        codice e nel diagramma, varianti, con che cosa non confonderlo), la \textbf{struttura} generale e un
        \textbf{esempio} completo, con il codice e il suo diagramma delle classi.
\end{itemize}

Il numero di ogni voce è un \textbf{collegamento}: nel PDF i rimandi come \textcolor{swblue!80!black}{\Lato\bfseries
1.6} portano alla voce citata; per chi stampa, l'indice dà le pagine. I colori distinguono i tipi di riquadro:
\textcolor{ncolore}{\textbf{verde petrolio}} per le voci di notazione, \textcolor{ccolore}{\textbf{blu}} per i casi
dal codice al diagramma, \textcolor{pcolore}{\textbf{viola}} per i pattern, \textcolor{acolore}{\textbf{arancione}}
per gli errori da evitare. L'etichetta \oltre{} segna ciò che nelle slide non c'è: costrutti di Java e notazioni UML
che il corso non mostra (mappe, matrici, classi generiche, \code{record}, \ldots), aggiunti per non trovarsi
impreparati se un esercizio li contiene.

Tutto il codice è stato compilato con Java 21, e gli esempi dei pattern sono stati eseguiti: gli output riportati
sono quelli reali. Per brevità negli esempi gli \code{import} di \code{java.util} sono sottintesi, e più classi stanno
nello stesso file.
